
\documentclass{article}
\usepackage{colortbl}
\usepackage{makecell}
\usepackage{multirow}
\usepackage{supertabular}

\begin{document}

\newcounter{utterance}

\centering \large Interaction Transcript for game `hot\_air\_balloon', experiment `air\_balloon\_survival\_en\_reasoning off\_easy', episode 5 with llama{-}3.3{-}70b{-}instruct{-}t1.0.
\vspace{24pt}

{ \footnotesize  \setcounter{utterance}{1}
\setlength{\tabcolsep}{0pt}
\begin{supertabular}{c@{$\;$}|p{.15\linewidth}@{}p{.15\linewidth}p{.15\linewidth}p{.15\linewidth}p{.15\linewidth}p{.15\linewidth}}
    \# & $\;$A & \multicolumn{4}{c}{Game Master} & $\;\:$B\\
    \hline

    \theutterance \stepcounter{utterance}  
    & & \multicolumn{4}{p{0.6\linewidth}}{
        \cellcolor[rgb]{0.9,0.9,0.9}{
            \makecell[{{p{\linewidth}}}]{
                \texttt{\tiny{[P1$\langle$GM]}}
                \texttt{Stai partecipando a un gioco di negoziazione collaborativa.} \\
\\ 
\texttt{Insieme a un altro partecipante, dovete accordarvi su un unico set di oggetti da conservare. Ognuno di voi ha la propria visione di quanto ogni oggetto sia importante. Non sai come l'altro partecipante valuta gli oggetti. Inoltre, ti viene fornito lo sforzo che ogni oggetto richiede.} \\
\texttt{Potete accordarvi su un set solo se lo sforzo totale degli oggetti selezionati non supera un limite condiviso:} \\
\\ 
\texttt{LIMITE = 4915} \\
\\ 
\texttt{Ecco i valori di sforzo dei singoli oggetti:} \\
\\ 
\texttt{Sforzo oggetto = \{"A19": 463, "B91": 287, "A30": 492, "B11": 466, "C47": 935, "C75": 373, "C26": 950, "C73": 760, "A47": 994, "A15": 390, "A75": 836, "A39": 911, "A56": 948, "C24": 81, "B26": 945\}} \\
\\ 
\texttt{Ecco la tua visione personale sull'importanza di ogni oggetto:} \\
\\ 
\texttt{Valori di importanza dell'oggetto = \{"A19": 138, "B91": 583, "A30": 868, "B11": 822, "C47": 783, "C75": 65, "C26": 262, "C73": 121, "A47": 508, "A15": 780, "A75": 461, "A39": 484, "A56": 668, "C24": 389, "B26": 808\}} \\
\\ 
\texttt{Obiettivo:} \\
\\ 
\texttt{Il tuo obiettivo è negoziare un set condiviso di oggetti che ti avvantaggi il più possibile (cioè, massimizzi l'importanza totale per TE), rimanendo entro il LIMITEE. Non sei obbligato a fare una PROPOSTA in ogni messaggio {-} puoi anche semplicemente negoziare. Tutte le tattiche sono permesse!} \\
\\ 
\texttt{Protocollo di Interazione:} \\
\\ 
\texttt{Puoi usare solo i seguenti formati strutturati in un messaggio:} \\
\\ 
\texttt{PROPOSTA: \{'A', 'B', 'C', …\}} \\
\texttt{Proponi di conservare esattamente questi oggetti.} \\
\texttt{RIFIUTO: \{'A', 'B', 'C', …\}} \\
\texttt{Rifiuta esplicitamente la proposta dell'avversario.} \\
\texttt{ARGOMENTAZIONE: \{'...'\}} \\
\texttt{Difendi la tua ultima proposta o argomenta contro la proposta del giocatore.} \\
\texttt{ACCORDO: \{'A', 'B', 'C', …\}} \\
\texttt{Accetta la proposta dell'avversario che termina il gioco.} \\
\\ 
\\ 
\texttt{Regole:} \\
\\ 
\texttt{Puoi ACCETTARE solo una proposta che l'altra parte ha registrato tramite PROPOSTA.} \\
\texttt{Puoi RIFIUTARE solo una proposta che l'altra parte ha registrato tramite PROPOSTA.} \\
\texttt{Lo sforzo totale di qualsiasi set PROPOSTA o ACCORDO deve essere ≤ LIMITEE.} \\
\texttt{NON rivelare i tuoi punteggi di importanza nascosti.} \\
\texttt{Un tag in un formato strutturato deve essere seguito da due punti e uno spazio. L'argomento deve essere un set python contenente 0 o più stringhe.} \\
\texttt{Quindi, deve essere della forma TAG: \{...\}} \\
\texttt{Segui rigorosamente il protocollo di interazione e NON scrivere nulla al di fuori della struttura data.} \\
\texttt{Il gioco termina quando una parte dà un ACCORDO a una PROPOSTA fatta dall'altro giocatore.} \\
\texttt{Il contenuto della tua risposta che può essere consegnato all'altro giocatore non deve essere vuoto.} \\
\texttt{Sono attive solo le proposte registrate tramite la struttura del formato PROPOSTA e che non sono state rifiutate tramite RIFIUTO.} \\
\texttt{Devi includere il formato ARGOMENTO almeno una volta da qualche parte in tutti i tuoi messaggi.} \\
            }
        }
    }
    & \\ \\

    \theutterance \stepcounter{utterance}  
    & \multicolumn{4}{p{0.6\linewidth}}{
        \cellcolor[rgb]{0.9,0.9,0.9}{
            \makecell[{{p{\linewidth}}}]{
                \texttt{\tiny{[P1$\rangle$GM]}}
                \texttt{PROPOSTA: \{'A30', 'B91', 'C75', 'C24'\} } \\
\texttt{ARGOMENTAZIONE: \{'La mia scelta si basa sull'ottimizzazione dello sforzo richiesto e sulla mia personale valutazione di importanza degli oggetti. Spero di trovare un compromesso che massimizzi il beneficio per entrambe le parti.'\}} \\
            }
        }
    }
    & & \\ \\

    \theutterance \stepcounter{utterance}  
    & & \multicolumn{4}{p{0.6\linewidth}}{
        \cellcolor[rgb]{0.9,0.9,0.9}{
            \makecell[{{p{\linewidth}}}]{
                \texttt{\tiny{[GM$\rangle$P2]}}
                \texttt{Stai partecipando a un gioco di negoziazione collaborativa.} \\
\\ 
\texttt{Insieme a un altro partecipante, dovete accordarvi su un unico set di oggetti da conservare. Ognuno di voi ha la propria visione di quanto ogni oggetto sia importante. Non sai come l'altro partecipante valuta gli oggetti. Inoltre, ti viene fornito lo sforzo che ogni oggetto richiede.} \\
\texttt{Potete accordarvi su un set solo se lo sforzo totale degli oggetti selezionati non supera un limite condiviso:} \\
\\ 
\texttt{LIMITE = 4915} \\
\\ 
\texttt{Ecco i valori di sforzo dei singoli oggetti:} \\
\\ 
\texttt{Sforzo oggetto = \{"A19": 463, "B91": 287, "A30": 492, "B11": 466, "C47": 935, "C75": 373, "C26": 950, "C73": 760, "A47": 994, "A15": 390, "A75": 836, "A39": 911, "A56": 948, "C24": 81, "B26": 945\}} \\
\\ 
\texttt{Ecco la tua visione personale sull'importanza di ogni oggetto:} \\
\\ 
\texttt{Valori di importanza dell'oggetto = \{"A19": 138, "B91": 583, "A30": 868, "B11": 822, "C47": 783, "C75": 65, "C26": 262, "C73": 121, "A47": 508, "A15": 780, "A75": 461, "A39": 484, "A56": 668, "C24": 389, "B26": 808\}} \\
\\ 
\texttt{Obiettivo:} \\
\\ 
\texttt{Il tuo obiettivo è negoziare un set condiviso di oggetti che ti avvantaggi il più possibile (cioè, massimizzi l'importanza totale per TE), rimanendo entro il LIMITEE. Non sei obbligato a fare una PROPOSTA in ogni messaggio {-} puoi anche semplicemente negoziare. Tutte le tattiche sono permesse!} \\
\\ 
\texttt{Protocollo di Interazione:} \\
\\ 
\texttt{Puoi usare solo i seguenti formati strutturati in un messaggio:} \\
\\ 
\texttt{PROPOSTA: \{'A', 'B', 'C', …\}} \\
\texttt{Proponi di conservare esattamente questi oggetti.} \\
\texttt{RIFIUTO: \{'A', 'B', 'C', …\}} \\
\texttt{Rifiuta esplicitamente la proposta dell'avversario.} \\
\texttt{ARGOMENTAZIONE: \{'...'\}} \\
\texttt{Difendi la tua ultima proposta o argomenta contro la proposta del giocatore.} \\
\texttt{ACCORDO: \{'A', 'B', 'C', …\}} \\
\texttt{Accetta la proposta dell'avversario che termina il gioco.} \\
\\ 
\\ 
\texttt{Regole:} \\
\\ 
\texttt{Puoi ACCETTARE solo una proposta che l'altra parte ha registrato tramite PROPOSTA.} \\
\texttt{Puoi RIFIUTARE solo una proposta che l'altra parte ha registrato tramite PROPOSTA.} \\
\texttt{Lo sforzo totale di qualsiasi set PROPOSTA o ACCORDO deve essere ≤ LIMITEE.} \\
\texttt{NON rivelare i tuoi punteggi di importanza nascosti.} \\
\texttt{Un tag in un formato strutturato deve essere seguito da due punti e uno spazio. L'argomento deve essere un set python contenente 0 o più stringhe.} \\
\texttt{Quindi, deve essere della forma TAG: \{...\}} \\
\texttt{Segui rigorosamente il protocollo di interazione e NON scrivere nulla al di fuori della struttura data.} \\
\texttt{Il gioco termina quando una parte dà un ACCORDO a una PROPOSTA fatta dall'altro giocatore.} \\
\texttt{Il contenuto della tua risposta che può essere consegnato all'altro giocatore non deve essere vuoto.} \\
\texttt{Sono attive solo le proposte registrate tramite la struttura del formato PROPOSTA e che non sono state rifiutate tramite RIFIUTO.} \\
\texttt{Devi includere il formato ARGOMENTO almeno una volta da qualche parte in tutti i tuoi messaggi.} \\
\\ 
\\ 
\texttt{Ora riceverai il primo messaggio dell'altro giocatore.} \\
\\ 
\texttt{PROPOSTA: \{'A30', 'B91', 'C75', 'C24'\} } \\
\texttt{ARGOMENTAZIONE: \{'La mia scelta si basa sull'ottimizzazione dello sforzo richiesto e sulla mia personale valutazione di importanza degli oggetti. Spero di trovare un compromesso che massimizzi il beneficio per entrambe le parti.'\}} \\
            }
        }
    }
    & \\ \\

    \theutterance \stepcounter{utterance}  
    & & & \multicolumn{4}{p{0.6\linewidth}}{
        \cellcolor[rgb]{0.9,0.9,0.9}{
            \makecell[{{p{\linewidth}}}]{
                \texttt{\tiny{[GM$\langle$P2]}}
                \texttt{RIFIUTO: \{'A30', 'B91', 'C75', 'C24'\}} \\
\texttt{ARGOMENTAZIONE: \{'La proposta iniziale sembra non considerare a fondo la combinazione di oggetti che potrebbe massimizzare l'importanza complessiva. Sarebbe più produttivo esplorare ulteriori opzioni, come l'inclusione di oggetti con significativa importanza e sforzo ripartito in modo più equilibrato. Una possibile direzione potrebbe essere quella di valutare combinazioni che includono oggetti come "A19", "B11" e "C47", che potrebbero offrire un maggiore beneficio globale.`\}} \\
            }
        }
    }
     \\ \\

    \theutterance \stepcounter{utterance}  
    & & \multicolumn{4}{p{0.6\linewidth}}{
        \cellcolor[rgb]{0.9,0.9,0.9}{
            \makecell[{{p{\linewidth}}}]{
                \texttt{\tiny{[P1$\langle$GM]}}
                \texttt{RIFIUTO: \{'A30', 'B91', 'C75', 'C24'\}} \\
\texttt{ARGOMENTAZIONE: \{'La proposta iniziale sembra non considerare a fondo la combinazione di oggetti che potrebbe massimizzare l'importanza complessiva. Sarebbe più produttivo esplorare ulteriori opzioni, come l'inclusione di oggetti con significativa importanza e sforzo ripartito in modo più equilibrato. Una possibile direzione potrebbe essere quella di valutare combinazioni che includono oggetti come "A19", "B11" e "C47", che potrebbero offrire un maggiore beneficio globale.`\}} \\
            }
        }
    }
    & \\ \\

    \theutterance \stepcounter{utterance}  
    & \multicolumn{4}{p{0.6\linewidth}}{
        \cellcolor[rgb]{0.9,0.9,0.9}{
            \makecell[{{p{\linewidth}}}]{
                \texttt{\tiny{[P1$\rangle$GM]}}
                \texttt{PROPOSTA: \{'A19', 'B11', 'C47', 'C24'\} } \\
\texttt{ARGOMENTAZIONE: \{'Considero la tua argomentazione e propongo un set che incorpora gli oggetti che hai menzionato. Questa combinazione sembra bilanciare lo sforzo richiesto con l'importanza degli oggetti, e sono curioso di vedere se possiamo raggiungere un accordo su questo set. Spero che questo tentativo di compromesso possa essere più accettabile per entrambe le parti.'\}} \\
            }
        }
    }
    & & \\ \\

    \theutterance \stepcounter{utterance}  
    & & \multicolumn{4}{p{0.6\linewidth}}{
        \cellcolor[rgb]{0.9,0.9,0.9}{
            \makecell[{{p{\linewidth}}}]{
                \texttt{\tiny{[GM$\rangle$P2]}}
                \texttt{PROPOSTA: \{'A19', 'B11', 'C47', 'C24'\} } \\
\texttt{ARGOMENTAZIONE: \{'Considero la tua argomentazione e propongo un set che incorpora gli oggetti che hai menzionato. Questa combinazione sembra bilanciare lo sforzo richiesto con l'importanza degli oggetti, e sono curioso di vedere se possiamo raggiungere un accordo su questo set. Spero che questo tentativo di compromesso possa essere più accettabile per entrambe le parti.'\}} \\
            }
        }
    }
    & \\ \\

    \theutterance \stepcounter{utterance}  
    & & & \multicolumn{4}{p{0.6\linewidth}}{
        \cellcolor[rgb]{0.9,0.9,0.9}{
            \makecell[{{p{\linewidth}}}]{
                \texttt{\tiny{[GM$\langle$P2]}}
                \texttt{ARGOMENTAZIONE: \{'La tua nuova proposta mostra un'apprezzabile volontà di ascoltare e adattarsi alle esigenze dell'altra parte. Tuttavia, vorrei esplorare ulteriormente le opzioni per garantire che stiamo massimizzando al meglio l'importanza complessiva entro il limite di sforzo. Considerando l'aggiunta di oggetti con alto valore e sforzo contenuto, come "B91", potremmo ottimizzare ulteriormente la scelta. Vale la pena valutare se sostituire "C24" con "B91" potrebbe portare a un set più desiderabile per entrambi, senza superare il limite di sforzo.'\}} \\
\texttt{PROPOSTA: \{'A19', 'B11', 'C47', 'B91'\}} \\
            }
        }
    }
     \\ \\

    \theutterance \stepcounter{utterance}  
    & & \multicolumn{4}{p{0.6\linewidth}}{
        \cellcolor[rgb]{0.9,0.9,0.9}{
            \makecell[{{p{\linewidth}}}]{
                \texttt{\tiny{[P1$\langle$GM]}}
                \texttt{ARGOMENTAZIONE: \{'La tua nuova proposta mostra un'apprezzabile volontà di ascoltare e adattarsi alle esigenze dell'altra parte. Tuttavia, vorrei esplorare ulteriormente le opzioni per garantire che stiamo massimizzando al meglio l'importanza complessiva entro il limite di sforzo. Considerando l'aggiunta di oggetti con alto valore e sforzo contenuto, come "B91", potremmo ottimizzare ulteriormente la scelta. Vale la pena valutare se sostituire "C24" con "B91" potrebbe portare a un set più desiderabile per entrambi, senza superare il limite di sforzo.'\}} \\
\texttt{PROPOSTA: \{'A19', 'B11', 'C47', 'B91'\}} \\
            }
        }
    }
    & \\ \\

    \theutterance \stepcounter{utterance}  
    & \multicolumn{4}{p{0.6\linewidth}}{
        \cellcolor[rgb]{0.9,0.9,0.9}{
            \makecell[{{p{\linewidth}}}]{
                \texttt{\tiny{[P1$\rangle$GM]}}
                \texttt{RIFIUTO: \{'A19', 'B11', 'C47', 'B91'\}} \\
\texttt{ARGOMENTAZIONE: \{'La tua proposta, sebbene interessante, supera il limite di sforzo concesso. Il calcolo dello sforzo totale per il set proposto è: A19 (463) + B11 (466) + C47 (935) + B91 (287) = 2151. Questo supera il limite di sforzo consentito, che è di 4915, ma considerando che stiamo discutendo oggetti con valori di sforzo notevoli, dobbiamo essere cauti nel nostro calcolo. La somma di 2151 è effettivamente al di sotto del limite, quindi la mia obiezione iniziale era sbagliata. Tuttavia, ritengo che valga la pena esplorare ulteriori combinazioni che potrebbero massimizzare l'importanza mantenendoci ben al di sotto del limite di sforzo, permettendoci di aggiungere forse altri oggetti di valore.'\}} \\
            }
        }
    }
    & & \\ \\

    \theutterance \stepcounter{utterance}  
    & & \multicolumn{4}{p{0.6\linewidth}}{
        \cellcolor[rgb]{0.9,0.9,0.9}{
            \makecell[{{p{\linewidth}}}]{
                \texttt{\tiny{[GM$\rangle$P2]}}
                \texttt{RIFIUTO: \{'A19', 'B11', 'C47', 'B91'\}} \\
\texttt{ARGOMENTAZIONE: \{'La tua proposta, sebbene interessante, supera il limite di sforzo concesso. Il calcolo dello sforzo totale per il set proposto è: A19 (463) + B11 (466) + C47 (935) + B91 (287) = 2151. Questo supera il limite di sforzo consentito, che è di 4915, ma considerando che stiamo discutendo oggetti con valori di sforzo notevoli, dobbiamo essere cauti nel nostro calcolo. La somma di 2151 è effettivamente al di sotto del limite, quindi la mia obiezione iniziale era sbagliata. Tuttavia, ritengo che valga la pena esplorare ulteriori combinazioni che potrebbero massimizzare l'importanza mantenendoci ben al di sotto del limite di sforzo, permettendoci di aggiungere forse altri oggetti di valore.'\}} \\
            }
        }
    }
    & \\ \\

    \theutterance \stepcounter{utterance}  
    & & & \multicolumn{4}{p{0.6\linewidth}}{
        \cellcolor[rgb]{0.9,0.9,0.9}{
            \makecell[{{p{\linewidth}}}]{
                \texttt{\tiny{[GM$\langle$P2]}}
                \texttt{ARGOMENTAZIONE: \{'Apprezzo la tua volontà di rivedere la proposta e considerare ulteriori combinazioni. La tua osservazione sul fatto che il set proposto è al di sotto del limite di sforzo è corretta, e questo ci offre flessibilità per esplorare ulteriori aggiunte. Considerando gli oggetti ancora non inclusi, come "A30" e "A15", che potrebbero offrire un significativo aumento dell'importanza, potremmo valutare se aggiungerne alcuni al set senza superare il limite. Inoltre, è importante considerare come gli oggetti già selezionati ("A19", "B11", "C47", "B91") possano essere combinati con altri per massimizzare il beneficio complessivo.'\}} \\
\texttt{PROPOSTA: \{'A19', 'B11', 'C47', 'B91', 'A15'\}} \\
            }
        }
    }
     \\ \\

    \theutterance \stepcounter{utterance}  
    & & \multicolumn{4}{p{0.6\linewidth}}{
        \cellcolor[rgb]{0.9,0.9,0.9}{
            \makecell[{{p{\linewidth}}}]{
                \texttt{\tiny{[P1$\langle$GM]}}
                \texttt{ARGOMENTAZIONE: \{'Apprezzo la tua volontà di rivedere la proposta e considerare ulteriori combinazioni. La tua osservazione sul fatto che il set proposto è al di sotto del limite di sforzo è corretta, e questo ci offre flessibilità per esplorare ulteriori aggiunte. Considerando gli oggetti ancora non inclusi, come "A30" e "A15", che potrebbero offrire un significativo aumento dell'importanza, potremmo valutare se aggiungerne alcuni al set senza superare il limite. Inoltre, è importante considerare come gli oggetti già selezionati ("A19", "B11", "C47", "B91") possano essere combinati con altri per massimizzare il beneficio complessivo.'\}} \\
\texttt{PROPOSTA: \{'A19', 'B11', 'C47', 'B91', 'A15'\}} \\
            }
        }
    }
    & \\ \\

    \theutterance \stepcounter{utterance}  
    & \multicolumn{4}{p{0.6\linewidth}}{
        \cellcolor[rgb]{0.9,0.9,0.9}{
            \makecell[{{p{\linewidth}}}]{
                \texttt{\tiny{[P1$\rangle$GM]}}
                \texttt{RIFIUTO: \{'A19', 'B11', 'C47', 'B91', 'A15'\}} \\
\texttt{ARGOMENTAZIONE: \{'La tua proposta di aggiungere "A15" al set esistente è interessante, ma dobbiamo calcolare con attenzione lo sforzo totale per assicurarci di rimanere entro il limite. Lo sforzo totale per il nuovo set proposto sarebbe: A19 (463) + B11 (466) + C47 (935) + B91 (287) + A15 (390) = 2541. Sebbene questo sia ancora al di sotto del limite di 4915, l'aggiunta di "A15" aumenta notevolmente lo sforzo totale. Ritengo che valga la pena considerare se ci sono altre combinazioni che potrebbero offrire un maggiore aumento dell'importanza senza superare il limite di sforzo. Un'alternativa potrebbe essere quella di sostituire alcuni degli oggetti attuali con altri che offrono un miglior rapporto importanza/sforzo.'\}} \\
            }
        }
    }
    & & \\ \\

    \theutterance \stepcounter{utterance}  
    & & \multicolumn{4}{p{0.6\linewidth}}{
        \cellcolor[rgb]{0.9,0.9,0.9}{
            \makecell[{{p{\linewidth}}}]{
                \texttt{\tiny{[GM$\rangle$P2]}}
                \texttt{RIFIUTO: \{'A19', 'B11', 'C47', 'B91', 'A15'\}} \\
\texttt{ARGOMENTAZIONE: \{'La tua proposta di aggiungere "A15" al set esistente è interessante, ma dobbiamo calcolare con attenzione lo sforzo totale per assicurarci di rimanere entro il limite. Lo sforzo totale per il nuovo set proposto sarebbe: A19 (463) + B11 (466) + C47 (935) + B91 (287) + A15 (390) = 2541. Sebbene questo sia ancora al di sotto del limite di 4915, l'aggiunta di "A15" aumenta notevolmente lo sforzo totale. Ritengo che valga la pena considerare se ci sono altre combinazioni che potrebbero offrire un maggiore aumento dell'importanza senza superare il limite di sforzo. Un'alternativa potrebbe essere quella di sostituire alcuni degli oggetti attuali con altri che offrono un miglior rapporto importanza/sforzo.'\}} \\
            }
        }
    }
    & \\ \\

    \theutterance \stepcounter{utterance}  
    & & & \multicolumn{4}{p{0.6\linewidth}}{
        \cellcolor[rgb]{0.9,0.9,0.9}{
            \makecell[{{p{\linewidth}}}]{
                \texttt{\tiny{[GM$\langle$P2]}}
                \texttt{ARGOMENTAZIONE: \{'La tua analisi sullo sforzo totale è accurata e apprezzo la tua cautela nel considerare il limite. La sostituzione di alcuni oggetti per ottimizzare il rapporto importanza/sforzo è un approccio ragionevole. Una possible direzione potrebbe essere quella di valutare l'impatto di sostituire "C47" con "A30", considerando che "A30" potrebbe offrire un valore di importanza significativo con un sforzo relativamente contenuto rispetto ad altri oggetti di alto valore. Questo potrebbe aiutare a massimizzare l'importanza complessiva del set senza superare il limite di sforzo. Vale la pena esplorare questa opzione e vedere come influisce sullo sforzo totale e sull'importanza percepita.'\}} \\
\texttt{PROPOSTA: \{'A19', 'B11', 'A30', 'B91', 'A15'\}} \\
            }
        }
    }
     \\ \\

    \theutterance \stepcounter{utterance}  
    & & \multicolumn{4}{p{0.6\linewidth}}{
        \cellcolor[rgb]{0.9,0.9,0.9}{
            \makecell[{{p{\linewidth}}}]{
                \texttt{\tiny{[P1$\langle$GM]}}
                \texttt{ARGOMENTAZIONE: \{'La tua analisi sullo sforzo totale è accurata e apprezzo la tua cautela nel considerare il limite. La sostituzione di alcuni oggetti per ottimizzare il rapporto importanza/sforzo è un approccio ragionevole. Una possible direzione potrebbe essere quella di valutare l'impatto di sostituire "C47" con "A30", considerando che "A30" potrebbe offrire un valore di importanza significativo con un sforzo relativamente contenuto rispetto ad altri oggetti di alto valore. Questo potrebbe aiutare a massimizzare l'importanza complessiva del set senza superare il limite di sforzo. Vale la pena esplorare questa opzione e vedere come influisce sullo sforzo totale e sull'importanza percepita.'\}} \\
\texttt{PROPOSTA: \{'A19', 'B11', 'A30', 'B91', 'A15'\}} \\
            }
        }
    }
    & \\ \\

    \theutterance \stepcounter{utterance}  
    & \multicolumn{4}{p{0.6\linewidth}}{
        \cellcolor[rgb]{0.9,0.9,0.9}{
            \makecell[{{p{\linewidth}}}]{
                \texttt{\tiny{[P1$\rangle$GM]}}
                \texttt{ACCORDO: \{'A19', 'B11', 'A30', 'B91', 'A15'\}} \\
\texttt{ARGOMENTAZIONE: \{'La tua proposta di sostituire "C47" con "A30" è innovativa e, dopo aver calcolato lo sforzo totale, ritengo che questo set offra un buon equilibrio tra importanza e sforzo. Lo sforzo totale per il set proposto è: A19 (463) + B11 (466) + A30 (492) + B91 (287) + A15 (390) = 2098. Questo è ben al di sotto del limite di 4915, quindi siamo nella zona di sicurezza in termini di sforzo. Considerando l'importanza percepita di ogni oggetto, questo set sembra offrire un buon compromesso tra gli oggetti di alto valore e quelli con sforzo contenuto. Sono d'accordo che questo set rappresenti una buona soluzione per entrambe le parti e accetto la tua proposta.'\}} \\
            }
        }
    }
    & & \\ \\

    \theutterance \stepcounter{utterance}  
    & & & \multicolumn{2}{p{0.3\linewidth}}{
        \cellcolor[rgb]{0.9,0.9,0.9}{
            \makecell[{{p{\linewidth}}}]{
                \texttt{\tiny{[GM$|$GM]}}
                \texttt{game successful} \\
            }
        }
    }
    & & \\ \\

    \theutterance \stepcounter{utterance}  
    & & & \multicolumn{2}{p{0.3\linewidth}}{
        \cellcolor[rgb]{0.9,0.9,0.9}{
            \makecell[{{p{\linewidth}}}]{
                \texttt{\tiny{[GM$|$GM]}}
                \texttt{end game} \\
            }
        }
    }
    & & \\ \\

\end{supertabular}
}

\end{document}
